\ExerciseSection

\begin{enumerate}
\def\labelenumi{\arabic{enumi})}
\item
  全序群 G（不必交换）称为阿基米德的，如果由 \(\geqslant e\) 的元素（G 的单位元）组成的集 I 满足 §2 的公理 \((\mathrm{GR}_{\mathrm{IV}})\) 。证明：全序群 G 同构于加法群 R 的一个子群，当且仅当 G 是阿基米德的（按 G 中 \textgreater e 的元素所成的集有或没有最小元素分成两种情形讨论，并利用 §2 的命题 1）。
\item
  设 G 是全序群（不必交换），且元素多于一个。则拓扑 \(\mathcal{C}_{0}(G)\) （第 I 章，§ 2，习题 5）与 G 的群结构相容。若 G 关于此拓扑是连通的，证明 G 同构于加法群 R（利用第 IV 章，§ 2 的习题 7 以及 § 2 的命题 2）。
\item
  设 G 是（不必交换的）全序群。若赋有拓扑 \(\mathcal{C}_{0}(G)\) 的 G 是局部紧的且非离散的，则 G 局部同构于 R，且 G 的单位连通分支是同构于 R 的开子群（利用第 IV 章，§ 2 的习题 6 以及 § 2 的命题 2）。
\item
  设 G 是满足下列条件的拓扑群：\((R_{I})\) G 连通。
\end{enumerate}

\((R_{\mathrm{II}})\) 补集 \(G^{*}\) （单位元 \(e\) 在 \(G\) 中的补集）不连通。

则存在 G 到 R 上的连续双射同态（换言之，G 代数同构于 R，且其拓扑细于 R 的拓扑）。证明分几步进行：

\begin{enumerate}
\def\labelenumi{\alph{enumi})}
\item
  设 \((\mathbf{U}_i)_{1 \leqslant i \leqslant n}\) 是 \(\mathbf{G}^*\) 分成 \(\mathbf{G}^* (n \geqslant 2)\) 中开集的一个划分。证明每个 \(\mathbf{U}_i\) 在 \(\mathbf{G}\) 中是开的，\(e\) 属于每个 \(\overline{\mathbf{U}}_i\) ，并且 \(\mathbf{G}\) 是豪斯多夫的。由此推出，\(\overline{\mathbf{U}}_i = \mathbf{U}_i \cup \{e\}\) 在 \(\mathbf{G}\) 中的闭包都是连通的（第 I 章，§ 11，习题 4）。
\item
  设 A 是 \(U_{i}\) 的一个连通分支。证明对每个指标 \(j \neq i\) ，有 \(A^{-1}\overline{U}_{j} = A^{-1}\) （注意 \(A^{-1}\overline{U}_{j}\) 连通且包含 \(A^{-1}\) ，并且 \(A^{-1}\overline{U}_{j} \subset G^{*}\) 只要 \(j \neq i\)）。
\item
  证明 n 必等于 2{[}对每个指标 i = 1, 2, \ldots, n，取 \(A_{i}\) 为 \(U_{i}\) 的一个分支；若 \(j \neq i\) ，则由 b) 有 \(A_{i}^{-1}A_{j} \subset A_{i}^{-1}\) ，\(A_{j}^{-1}A_{i} \subset A_{j}^{-1} \subset U_{j}^{-1}\) ，从而 \(A_{i}^{-1} \subset U_{j}\){]}。由此推出：\(G^{*}\) 恰有两个分支 A 与 B，并且 \(B = A^{-1}\) ，\(\overline{A} = \complement(A^{-1})\) 。
\item
  关系 \(yx^{-1} \in \overline{\mathbf{A}}\) 是一个序关系，它使 \(\mathbf{G}\) 成为全序群（证明 \(\overline{\mathbf{A}}^2 = \overline{\mathbf{A}}\) ，并且 \(x\overline{\mathbf{A}}x^{-1} = \overline{\mathbf{A}}\) 对所有 \(x \in \mathbf{G}\) 成立）。
\item
  拓扑 \(\mathfrak{C}_0(\mathbf{G})\) 粗于给定的拓扑 \(\mathfrak{C}\) （\(\mathbf{G}\) 上的）（注意 \(\mathbf{A}\) 关于 \(\mathfrak{C}\) 是开的）。
\end{enumerate}

利用上面的习题 2 以及第 1 章，§ 11，第 2 小节，命题 4 完成证明。{[}参见第 VI 章，§ 1，习题 12 b)。{]}

\begin{enumerate}
\def\labelenumi{\alph{enumi})}
\setcounter{enumi}{5}
\tightlist
\item
  证明：若再假设 G 是局部紧的或局部连通的，则 G 同构于 R。
\end{enumerate}

* 5) 在 \(\mathbf{R}\) 上给出一个与群结构相容的拓扑的例子，使 \(\mathbf{R}\) 关于它是连通、局部紧且局部连通的，并且使 \(\{o\}\) 在 \(\mathbf{R}\) 中的补集是连通的（利用 \(\mathbf{R}\) 与 \(\mathbf{R}^n\) 代数同构这一事实）。 *

\begin{enumerate}
\def\labelenumi{\arabic{enumi})}
\setcounter{enumi}{5}
\tightlist
\item
  设 G 是满足下列条件的拓扑群：\((\mathrm{LR}_{\mathbf{I}})\) G 是豪斯多夫的且局部连通。
\end{enumerate}

\((\mathrm{LR}_{\mathbf{II}})\) 存在一个连通邻域 \(\mathbf{U}\) （单位元 \(e\) 在 \(\mathbf{G}\) 中的邻域），使得 \(e\) 在 \(\mathbf{U}\) 中的补集不连通。证明：在这些条件下，\(\mathbf{G}\) 局部同构于 \(\mathbf{R}\) 。

{[}先证明：若 \(\mathbf{V}\) 是 \(e\) 在 \(\mathbf{G}\) 中的含于 \(\mathbf{U}\) 的任一连通邻域，则 \(\mathbf{V} \cap \complement\{e\}\) 不连通，\(e\) 属于 \(\mathbf{V} \cap \complement\{e\}\) 的每个分支的闭包，并且 \(\mathbf{V} \cap \complement\{e\}\) 恰有两个分支，其推理方式与习题 4 中的相同。然后取充分小的连通邻域 \(\mathbf{V}\) ，并在 \(\mathbf{V}\) 上定义一个全序，使得 \(\mathfrak{C}_0(\mathbf{V})\) 粗于 \(\mathbf{V}\) 上由 \(\mathbf{G}\) 的拓扑诱导的拓扑，并且使得 § 2 的命题 2 可应用于集 \(\mathbf{E}\) ，其元素 \(\geqslant e\) 且属于 \(\mathbf{V}\) 。{]}

\begin{enumerate}
\def\labelenumi{\arabic{enumi})}
\setcounter{enumi}{6}
\tightlist
\item
  设 G 是具有可数开集基的非离散局部紧群，并设 \(V_{0}\) 是 G 的单位元 e 的紧邻域，它除 \(\{e\}\) 外不含 G 的任何子群。
\end{enumerate}

\begin{enumerate}
\def\labelenumi{\alph{enumi})}
\tightlist
\item
  设 \((x_{n})\) 是 \(\mathbf{V}_{0}\) 中的收敛于 \(e\) 的点的序列。对每个 \(n\) ，存在整数 \(p(n) > 0\) ，使得 \(x_{n}^{k} \in \mathbf{V}_{0}\) ，只要 \(k \leqslant p(n)\) ，且 \(x_{n}^{p(n)+1} \notin \mathbf{V}_{0}\) ，并且我们有 \(\lim_{n \to \infty} p(n) = +\infty\) 。证明：存在子序列 \((x_{k_{n}})\) （\((x_{n})\) 的），使得对每个有理数 \(r\) （\(|r| \leqslant 1\)），序列 \((x_{k_{n}}^{[rp(k_{n})]})\) 有极限 \(f(r)\) （在 \(\mathbf{V}_{0}\) 中；利用对角线方法）；并且若 \(r, r'\) 是满足 \(|r| \leqslant 1\) ，\(|r'| \leqslant 1\) 与 \(|r + r'| \leqslant 1\) 的有理数，则我们有
\end{enumerate}

\[
f (r + r ^ {\prime}) = f (r) f (r ^ {\prime}).
\]

\begin{enumerate}
\def\labelenumi{\alph{enumi})}
\setcounter{enumi}{1}
\item
  证明 \(f\) 在 \(o\) 的一个邻域（\(\mathbf{Q}\) 中）内连续。{[}用反证法：若存在有理数序列 \((r_n)\) 趋于 \(o\) ，且使 \(f(r_n)\) 趋于元素 \(y \neq e\) （在 \(V_0\) 中），试证应有 \(y^m \in V_0\) ，对一切整数 \(m\) 。{]}
\item
  由 b) 推出：存在 \(\mathbf{R}\) 到 \(\mathbf{G}\) 中的非平凡连续同态（利用 §1 的命题 6）。
\item
  设 H 是 G 的闭正规子群，H 异于 G 本身，并假设在 G/H 中存在单位元的一个邻域，它不含任何非平凡子群。证明：存在连续同态 \(f: R \rightarrow G\) ，使得复合映射
\end{enumerate}

\[
\mathbf {R} \xrightarrow {f} \mathbf {G} \to \mathbf {G} / \mathbf {H}
\]

是非平凡的。

\begin{enumerate}
\def\labelenumi{\arabic{enumi})}
\setcounter{enumi}{7}
\tightlist
\item
  设 G 是拓扑群，H 是 G 的闭子群，含于 G 的中心，并且存在连续满射同态 \(\varphi: \mathbf{R} \to \mathbf{G} / \mathbf{H}\) 。证明 \(\mathbf{G}\) 是交换的{[}若 \(a_0\) 是 \(\mathbf{G}\) 中不属于 \(\mathbf{H}\) 的元素，注意存在元素序列 \((a_n)\) （\(\mathbf{G}\) 中的）与元素序列 \((c_n)\) （\(\mathbf{H}\) 中的），使得 \(a_n = c_n a_{n+1}\) ，并且 \(\mathbf{G}\) 的子群——由 \(\mathbf{H}\) 与诸 \(a_n\) 生成——在 \(\mathbf{G}\) 中稠密{]}。
\end{enumerate}

